\documentclass[11pt]{article}
\usepackage[margin=1in]{geometry}
\usepackage[T1]{fontenc}
\usepackage[utf8]{inputenc}
\usepackage{lmodern}
\usepackage{booktabs,longtable,graphicx,microtype}
\usepackage{setspace,lineno,hyperref}
\hypersetup{hidelinks}
\doublespacing
\linenumbers

\title{Supplementary material\\
Metropolitan scope amplifies measured gender disparities in stochastic urban accessibility}
\author{Andy Rafael Domínguez-Monterroza}
\date{}

\begin{document}
\maketitle

\section{Supplementary methods}

\subsection{Interpretation of the stochastic-accessibility estimand}
The origin--destination Markov chain represents aggregate survey-weighted population
flows. A first-passage step is a dimensionless transition in this fitted flow system and
must not be interpreted as a minute of travel or as an observed individual's itinerary.
For every purpose, residential distributions are conditioned on starting outside the
corresponding target set in the primary specification.

\subsection{Regularization sensitivity}
Group-specific rows are shrunk toward the population kernel using
$\omega_{ig}=\mathrm{ESS}_{ig}/(\mathrm{ESS}_{ig}+\tau)$. The primary analysis uses
$\tau=20$. We re-estimated both territorial systems at $\tau\in\{0,10,20,50\}$.
At $\tau=0$, origins with positive group-specific support use the empirical row, whereas
origins without such support retain the population row so that the transition kernel
remains defined. Table~\ref{tab:s1_tau} reports the total gap, its paired scope contrast,
the corresponding residential and kernel contrasts, and the normalized-gap contrast.
It also retains the origin-level magnitude and rank diagnostics relative to $\tau=20$.
These are point-estimate sensitivity calculations; inferential intervals in the main
analysis correspond to the primary $\tau=20$ specification.

\begin{table}
\caption{Sensitivity of group FPT estimates to the regularization parameter.}
\label{tab:s1_tau}
\begin{tabular}{rrrr}
\toprule
Tau & Max. absolute change & Min. Spearman & Median Spearman \\
\midrule
10.000 & 0.078 & 0.949 & 0.989 \\
20.000 & 0.000 & 1.000 & 1.000 \\
50.000 & 0.103 & 0.971 & 0.990 \\
\bottomrule
\end{tabular}
\end{table}


\subsection{Specification grid}
The robustness grid crosses two reference populations (all survey-valid observations or
only observations with valid categories for the analysed dimension), two target
definitions (core 50\% and extended 80\% of supported weighted arrivals), and two
residential conventions (outside-target conditioning or all residences). Across three
purposes this yields 24 sex-gap specifications. These checks assess specification
dependence; they were not used to select the primary result.

\begin{table}
\caption{Range of regional female-minus-male decomposition estimates over the 12 effective full-population robustness specifications.}
\label{tab:s2_ranges}
\begin{tabular}{llll}
\toprule
Purpose & Total gap range & Residential range & Kernel range \\
\midrule
Education & [0.090, 0.827] & [-0.013, 0.030] & [0.077, 0.840] \\
Employment & [0.127, 0.999] & [-0.043, -0.020] & [0.147, 1.031] \\
Health & [0.493, 0.959] & [-0.046, -0.031] & [0.532, 0.997] \\
\bottomrule
\end{tabular}
\end{table}

\begin{longtable}{llllrrr}
\caption{Female-minus-male decomposition for the complete prespecified robustness grid.} \label{tab:s3_grid} \\
\toprule
Reference & Target & Conditioning & Purpose & Total gap & Residential & Kernel \\
\midrule
\endfirsthead
\caption[]{Female-minus-male decomposition for the complete prespecified robustness grid.} \\
\toprule
Reference & Target & Conditioning & Purpose & Total gap & Residential & Kernel \\
\midrule
\endhead
\midrule
\multicolumn{7}{r}{Continued on next page} \\
\midrule
\endfoot
\bottomrule
\endlastfoot
Full population & Core 50\% & Outside target & Employment & 0.999 & -0.033 & 1.031 \\
Full population & Core 50\% & All residences & Employment & 0.754 & -0.043 & 0.797 \\
Full population & Extended 80\% & Outside target & Employment & 0.358 & -0.020 & 0.378 \\
Full population & Extended 80\% & All residences & Employment & 0.127 & -0.020 & 0.147 \\
Full population & Core 50\% & Outside target & Education & 0.827 & -0.013 & 0.840 \\
Full population & Core 50\% & All residences & Education & 0.592 & 0.030 & 0.562 \\
Full population & Extended 80\% & Outside target & Education & 0.239 & -0.012 & 0.251 \\
Full population & Extended 80\% & All residences & Education & 0.090 & 0.013 & 0.077 \\
Full population & Core 50\% & Outside target & Health & 0.705 & -0.031 & 0.736 \\
Full population & Core 50\% & All residences & Health & 0.570 & -0.046 & 0.616 \\
Full population & Extended 80\% & Outside target & Health & 0.959 & -0.038 & 0.997 \\
Full population & Extended 80\% & All residences & Health & 0.493 & -0.039 & 0.532 \\
Valid sex categories & Core 50\% & Outside target & Employment & 0.999 & -0.033 & 1.031 \\
Valid sex categories & Core 50\% & All residences & Employment & 0.754 & -0.043 & 0.797 \\
Valid sex categories & Extended 80\% & Outside target & Employment & 0.358 & -0.020 & 0.378 \\
Valid sex categories & Extended 80\% & All residences & Employment & 0.127 & -0.020 & 0.147 \\
Valid sex categories & Core 50\% & Outside target & Education & 0.827 & -0.013 & 0.840 \\
Valid sex categories & Core 50\% & All residences & Education & 0.592 & 0.030 & 0.562 \\
Valid sex categories & Extended 80\% & Outside target & Education & 0.239 & -0.012 & 0.251 \\
Valid sex categories & Extended 80\% & All residences & Education & 0.090 & 0.013 & 0.077 \\
Valid sex categories & Core 50\% & Outside target & Health & 0.705 & -0.031 & 0.736 \\
Valid sex categories & Core 50\% & All residences & Health & 0.570 & -0.046 & 0.616 \\
Valid sex categories & Extended 80\% & Outside target & Health & 0.959 & -0.038 & 0.997 \\
Valid sex categories & Extended 80\% & All residences & Health & 0.493 & -0.039 & 0.532 \\
\end{longtable}


\subsection{Multiscale territorial construction}
External municipalities are ranked by total bidirectional survey-weighted flow with
Bogotá using the complete population. Nested functional scopes include the minimum
ranked municipalities required to reach at least 50\% and 80\% of that interface flow.
Together with Bogotá D.C. and the complete regional network, the realized interface
shares are 0, 0.564, 0.824, and 1. Within every scope, the largest directed strongly
connected component, population kernel, purpose targets, group kernels, and residential
distributions are reconstructed. The full-region endpoint retains all 141 nodes in the
validated regional kernel, including one node without a municipality label.

\subsection{Bootstrap inference}
The final regional analysis uses 500 joint Poisson(1) household-cluster bootstrap
replicates. The same multiplier is assigned to every record from a household and is
propagated jointly to person and trip weights. Simultaneous intervals use a bootstrap
max-$t$ critical value over the three destination purposes within each effect family.
The paired scope analysis likewise uses 500 shared household multipliers in the regional
and district systems. The earlier 50-replicate specification grid and 200-replicate
standalone district analysis are robustness diagnostics rather than the final paired
inference.

\section{Supplementary results}

\subsection{Regional inference}
\begin{table}
\caption{Final Bogotá-Region household-bootstrap inference.}
\label{tab:s4_region}
\begin{tabular}{llrrrrrr}
\toprule
Purpose & Effect & Estimate & Bootstrap SD & Pointwise lower & Pointwise upper & Simultaneous lower & Simultaneous upper \\
\midrule
Education & Total female-minus-male gap & 0.827 & 0.168 & 0.547 & 1.194 & 0.451 & 1.204 \\
Education & Residential component & -0.013 & 0.024 & -0.054 & 0.037 & -0.067 & 0.042 \\
Education & Mobility-kernel component & 0.840 & 0.165 & 0.581 & 1.206 & 0.467 & 1.213 \\
Employment & Total female-minus-male gap & 0.999 & 0.164 & 0.718 & 1.370 & 0.632 & 1.366 \\
Employment & Residential component & -0.033 & 0.024 & -0.074 & 0.018 & -0.088 & 0.023 \\
Employment & Mobility-kernel component & 1.031 & 0.160 & 0.758 & 1.375 & 0.671 & 1.392 \\
Health & Total female-minus-male gap & 0.705 & 0.203 & 0.363 & 1.109 & 0.252 & 1.159 \\
Health & Residential component & -0.031 & 0.029 & -0.085 & 0.028 & -0.096 & 0.035 \\
Health & Mobility-kernel component & 0.736 & 0.196 & 0.411 & 1.133 & 0.294 & 1.178 \\
\bottomrule
\end{tabular}
\end{table}


\subsection{District support and target reconstruction}
The Bogotá D.C.-internal analysis retains the largest directed strongly connected
component. Purpose-specific targets are reconstructed within the district scope rather
than inherited from the regional system.

\begin{table}
\caption{Support and strongly connected component diagnostics for the Bogotá D.C.-internal network.}
\label{tab:s5_support}
\begin{tabular}{lr}
\toprule
Metric & Value \\
\midrule
City utam in zoning & 116.000 \\
Largest scc nodes & 114.000 \\
Internal trip rows before scc & 76168.000 \\
Internal trip rows in scc & 76168.000 \\
Trip weight scc coverage & 1.000 \\
City person rows before scc & 48989.000 \\
City person rows in scc & 48373.000 \\
Person weight scc coverage & 0.998 \\
\bottomrule
\end{tabular}
\end{table}

\begin{table}
\caption{Scope-specific Bogotá D.C.-internal opportunity target sets.}
\label{tab:s6_targets}
\begin{tabular}{llrr}
\toprule
Purpose & Definition & Target nodes & Supported arrival share \\
\midrule
Employment & core\_50 & 24 & 0.500 \\
Education & core\_50 & 24 & 0.512 \\
Health & core\_50 & 17 & 0.508 \\
\bottomrule
\end{tabular}
\end{table}


\subsection{Paired territorial-scope contrasts}
Positive step-valued contrasts indicate a larger female-minus-male gap in the regional
system. Relative attenuation is reported in percentage points. The scope comparison is
descriptive and includes both network truncation and scope-specific opportunity
geography.

\begin{table}
\caption{Paired Bogotá-Region minus Bogotá D.C.-internal scope contrasts.}
\label{tab:s7_paired}
\begin{tabular}{lllrrrr}
\toprule
Purpose & Effect & Unit & Estimate & Bootstrap SD & Simultaneous lower & Simultaneous upper \\
\midrule
Employment & Total gap & FPT steps & 0.520 & 0.153 & 0.178 & 0.861 \\
Employment & Residential component & FPT steps & -0.022 & 0.025 & -0.078 & 0.034 \\
Employment & Mobility-kernel component & FPT steps & 0.542 & 0.146 & 0.217 & 0.866 \\
Employment & Relative attenuation & Percent & 52.035 & 8.777 & 31.831 & 72.240 \\
Education & Total gap & FPT steps & 0.429 & 0.160 & 0.070 & 0.787 \\
Education & Residential component & FPT steps & -0.022 & 0.024 & -0.076 & 0.031 \\
Education & Mobility-kernel component & FPT steps & 0.451 & 0.155 & 0.108 & 0.794 \\
Education & Relative attenuation & Percent & 51.821 & 11.178 & 26.090 & 77.552 \\
Health & Total gap & FPT steps & 0.465 & 0.169 & 0.087 & 0.843 \\
Health & Residential component & FPT steps & -0.030 & 0.028 & -0.091 & 0.032 \\
Health & Mobility-kernel component & FPT steps & 0.495 & 0.162 & 0.136 & 0.853 \\
Health & Relative attenuation & Percent & 65.911 & 15.801 & 29.537 & 102.284 \\
\bottomrule
\end{tabular}
\end{table}


\subsection{Multiscale boundary response}
The total female-minus-male gap is regressed on realized interface-flow coverage within
each bootstrap replicate. Max-$t$ simultaneous intervals are formed across the three
purposes separately for each decomposition component. Table~\ref{tab:s8_multiscale}
reports the fitted slopes. Positive simultaneous intervals support an overall increasing
response, but do not imply that every successive scope estimate is larger.

\begin{table}[ht]
\centering
\caption{Multiscale boundary-response slopes per unit of Bogotá--external interface-flow coverage.}
\label{tab:s8_multiscale}
\begin{tabular}{lrrr}
\toprule
Purpose & Total gap & Residential & Mobility kernel \\
\midrule
Employment & 0.451 [0.161, 0.741] & -0.018 [-0.058, 0.023] & 0.469 [0.175, 0.762] \\
Education  & 0.320 [0.024, 0.616] & -0.018 [-0.057, 0.021] & 0.338 [0.035, 0.641] \\
Health     & 0.391 [0.052, 0.730] & -0.026 [-0.073, 0.022] & 0.417 [0.078, 0.755] \\
\bottomrule
\end{tabular}
\end{table}

Point-estimate monotonicity holds for employment but not for education or health.
The corresponding bootstrap probabilities of a nondecreasing total-gap sequence are
0.756, 0.110, and 0.212. The main inference is therefore the positive fitted response,
not strict local monotonicity.

\subsection{Socioeconomic-stratum diagnostic}
\begin{table}
\caption{Association between socioeconomic stratum and observed mean FPT.}
\label{tab:s9_stratum}
\begin{tabular}{lrrrr}
\toprule
Purpose & Spearman rho & Bootstrap SD & Lower 95\% CI & Upper 95\% CI \\
\midrule
Education & 0.886 & 0.141 & 0.314 & 0.886 \\
Employment & -1.000 & 0.046 & -1.000 & -0.829 \\
Health & -0.771 & 0.120 & -0.829 & -0.429 \\
\bottomrule
\end{tabular}
\end{table}


\section{Reproducibility and interpretation boundaries}
All supplementary tables are generated from passed analytical gates without
re-estimating the models. The analysis is cross-sectional and descriptive. Survey
weights and clustered bootstrap inference propagate sampling uncertainty but do not
remove recall error, nonresponse bias, unobserved modal and temporal heterogeneity, or
the limited gender categories in the processed survey.

\end{document}
